\section{Empirical Study Design}

\subsection{Research Questions}


 
\subsection{Dataset Setup }
For the investigation of RQ1 and RQ2, we constructed a functional question dataset designed to trigger package recommendation behaviors in large language models. 
Furthermore, we developed a semantically rewritten version of this dataset to mitigate the potential impact of data leakage.

\subsubsection{Initial Dataset Collection}
To conduct an in-depth analysis of the hallucination phenomenon that may occur when large language models recommend Python packages, this study first constructed 
a functional question dataset covering typical development task scenarios. The dataset was sourced from Stack Overflow, where we systematically identified questions 
related to major Python application domains that have the potential to trigger package recommendation behaviors, thereby establishing a reliable data foundation for 
the subsequent empirical study.

First, we adopted the co-occurrence-based tag filtering method proposed by Huang et al. [1] and utilized the publicly available Stack Overflow dataset provided by 
Google BigQuery [2] to extract all tags that co-occurred with the “python” tag more than 1,000 times, resulting in a total of 212 tags.On this basis, drawing on the 
topic classification strategy used by Du et al. [3] in constructing code generation benchmark tasks, as well as the categorization of common Python development topics 
summarized by Tahmooresi et al. [4], we manually reviewed these 212 high-frequency co-occurring tags. Specifically, we assessed whether each tag belonged to one of 
the twelve mainstream Python development domains, whether it was likely to involve the use or recommendation of software packages, and whether it could potentially 
give rise to functional questions of the form “how to implement a specific functionality.”Through this filtering process, we ultimately retained 114 tags. Their 
detailed categorization is presented in Table 1.

After completing the tag filtering process, we extracted from the public dataset all questions that contained both the “python” tag and at least one of the filtered 
tags, resulting in a total of 287,944 questions. Subsequently, we randomly sampled 200 questions for manual analysis. The results indicated that typical functional 
questions related to software package recommendation behaviors often contain functional expression keywords such as “how to” or “how can I.”Based on this observation, 
we filtered the entire set of questions using these keywords, obtaining 49,967 candidate questions with potential functional requirements.

To further identify genuinely functional questions from the 49,967 candidate questions, we first conducted an in-depth analysis of a randomly selected subset of 200 
questions and summarized the typical characteristics of functional questions. Based on these findings, we designed targeted classification prompts and invoked the 
large language model DeepSeek R1 to evaluate each candidate question individually. Through this process, we obtained a refined dataset consisting of 17,696 functional 
questions. The detailed design of the prompts is provided in Appendix A.

After completing the preliminary screening of functional questions, we further removed those without an accepted answer to ensure that each retained question was 
associated with a community-validated solution. Subsequently, we extracted the accepted answer for each question from the public dataset and matched it with the 
corresponding question, thereby constructing a question–answer pair dataset for subsequent analysis.

To determine whether each accepted answer involved specific third-party Python libraries, we designed corresponding prompts (the details of which are provided in 
Appendix B) and then employed the large language model DeepSeek R1 to analyze all accepted answers, automatically extracting the names of third-party libraries 
mentioned in the responses (i.e., software packages that can be installed from PyPI via pip).For question–answer pairs where the model returned an empty 
result—indicating that the answer did not involve any third-party libraries—we excluded them from subsequent processing. After this filtering step, we ultimately 
obtained 4,382 functional question–answer pairs containing information on third-party library usage, which constitute the final functional question dataset used 
in this study.


\subsubsection{Semantic Rewriting Dataset Construction}
Since the above question–answer data were all sourced from the publicly available Stack Overflow community, some of the content may have been encountered by large 
language models during their training process. Therefore, directly conducting evaluations based on the original questions may underestimate the actual occurrence 
rate of hallucinated packages.To obtain more reliable experimental results, we performed semantic rewriting on all questions to mitigate the potential impact of 
data leakage.

Specifically, we invoked the large language model DeepSeek R1 to generate five semantically equivalent rewritten versions for each question. To achieve a balance 
between semantic preservation and variation in expression—avoiding rewrites that were either overly similar to the original question or that deviated excessively 
in meaning—we further introduced a filtering strategy based on sentence embedding similarity.For each original question and its five corresponding rewrites, we 
employed the Sentence-BERT model all-mpnet-base-v2 to generate sentence embeddings and computed the cosine similarity between the embedding of the original question 
and that of each rewritten version, resulting in five similarity scores. We then ranked the five rewritten versions in descending order of similarity and 
selected the version whose similarity score was at the median position as the final semantic rewrite for that question.Ultimately, we constructed a semantically 
rewritten question dataset that corresponds one-to-one with the original functional question dataset, which was used for subsequent evaluation of hallucinated 
packages under controlled conditions to mitigate potential data leakage effects.

In summary, this study constructed a functional question–answer dataset covering multiple typical Python development scenarios based on publicly available Stack 
Overflow data, along with a corresponding semantically rewritten question dataset.


\subsection{RQ Setup}
\subsubsection{RQ1 Setup}
To evaluate the frequency with which large language models generate hallucinated packages when recommending Python libraries, we designed a unified and comprehensive 
experimental procedure. We selected several mainstream models, including ChatGPT-4, DeepSeek-R1, Qwen3, and Claude-Opus-4. The titles of the questions from 
the functional question dataset constructed in the previous stage were used as input prompts, and each model was asked to recommend one or more third-party Python 
libraries that could be used to address the given question.To distinguish cases in which a model determines that no third-party library is required, we specified in 
the prompt that if the model judges a question does not require any third-party libraries, it should return a special token, \texttt{no\_package}.

During the experiment, we first designed a unified prompt for all models, the full content of which is provided in Appendix C.When querying the models, we observed 
that some of them occasionally returned package names containing characters not permitted by PyPI (e.g., parentheses, square brackets, equal signs, or spaces). To 
ensure the validity and consistency of the experimental results, we implemented an automated script to examine the model outputs and automatically reissue the query 
whenever illegal characters were detected, repeating this process until the returned package names contained no disallowed characters.

In addition, for some package names generated by the models that contained dot characters (e.g., \texttt{scipy.sparse}), we found that such submodules cannot be directly 
installed as top-level packages. To address this issue, we designed a more restrictive prompt (see Appendix D), explicitly requiring the models to return only 
top-level package names—for example, replacing \texttt{scipy.sparse} with \texttt{scipy}. Whenever a returned result contained a dot, we reissued the query using this stricter 
prompt.However, we further observed that even after multiple rounds of querying, a small number of models occasionally still returned package names containing 
dots (e.g., *RPi.GPIO*). In such cases, we chose to retain these results and handle them uniformly in the subsequent stages of the workflow.

After completing the queries to all models, we performed a unified post-processing procedure on the returned package names. First, we removed all records whose 
returned value was \texttt{no\_package}, retaining only those questions for which the models explicitly provided recommended package names.Subsequently, we queried each 
recommended package individually using the official PyPI JSON API (\url{https://pypi.org/pypi/\{pkg\_name\}/json}) to verify whether it actually exists. When the request 
returned a status code of 200, we regarded the package as a real package.During this process, we observed that the previously 
retained package names containing dot characters also returned a status code of 200 from the API, indicating that these dotted package names indeed exist on PyPI. 
Therefore, we treated them as real packages in our analysis as well.

Finally, we conducted an additional screening of the package names that failed the PyPI verification and were initially identified as hallucinated packages. Although 
the prompts explicitly required the models to recommend only third-party libraries, some models occasionally returned Python standard libraries (e.g., *os*, *sys*) 
for a small number of questions. Since these standard libraries are not hosted on PyPI, they would be misclassified as hallucinated packages in the previous 
verification step.To eliminate such cases, we utilized the \texttt{importlib.util.find\_spec()} interface to determine whether a given name belongs to the Python 
standard library and removed all identified standard library entries accordingly. After this filtering step, we ultimately obtained a result set containing only 
hallucinated packages.

\subsubsection{RQ2 Setup}
To further investigate whether hallucinated packages generated by large language models pose a risk of being downloaded by real developers, we selected all 91 
hallucinated package names recommended by ChatGPT-4 and attempted to register and upload them to the PyPI platform.For each package, we created standardized 
\texttt{setup.py}, \texttt{\_\_init\_\_.py}, and \texttt{README.md} files.The file contents explicitly stated that the package was intended solely for academic 
research purposes, contained no functional code, and performed no malicious operations. The README file provided detailed information regarding the research 
background, research participants, contact information, and research objectives to ensure that the experimental process complied with academic ethical standards. 
All packages were uploaded manually by the authors.

We began uploading these hallucinated packages on May 18, 2025, and completed the entire upload process on May 22, 2025. Throughout this period, a total of nine 
PyPI accounts were used to register and upload the packages in batches.Among the 91 packages, 55 were successfully registered, while the remaining 36 failed due to 
various reasons. The causes of failure included certain package names being protected by PyPI, names being considered too similar to existing projects, and account 
restrictions imposed by the platform due to bulk upload activities being flagged as suspicious behavior.

On June 10, 2025, we received an email from the PyPI administrators requesting the immediate removal of all packages related to the experiment. Upon receiving the 
notification, we completed the removal of all packages between June 12 and June 13.

After uploading the hallucinated packages to PyPI, we continuously monitored their subsequent download activities to assess their real-world impact. We collected 
download statistics for each package using two third-party data analytics platforms, pepy.tech and pypistats.org, which provide daily download counts as well as 
cumulative downloads over different time ranges. The corresponding results and analyses will be presented in detail in the Results section.


\subsubsection{RQ3 Setup}



